
\chapter*{Resumen}
\label{chapter:0-Resumen}

La inteligencia artificial generativa ha transformado diversos ámbitos y áreas. Un campo particularmente impactado es la educación superior. En búsqueda de un aprendizaje más eficiente y accesible, los estudiantes recurren y consultan cada vez más a grandes modelos de lenguaje (LLMs) y otras herramientas de inteligencia artificial generativa para complementar sus estudios y resolver dudas académicas. Sin embargo, los modelos comerciales funcionan como cajas negras y su uso conlleva ciertos riesgos, por lo que surge la necesidad de proveer a los estudiantes de la FING una solución personalizada que mitigue esos riesgos.

En este documento se presenta el diseño, la implementación y la evaluación de un asistente de inteligencia artificial generativa orientado a estudiantes de la Facultad de Ingeniería (FING) de la Universidad de la República (UdelaR). El asistente implementado de inteligencia artificial generativa se basa en la arquitectura de Generación Aumentada por Recuperación (RAG) y se utilizan como corpus de conocimiento los videos a las clases publicadas en OpenFing. Para ello, se desarrolló un pipeline automatizado con el objetivo de transcribir mediante herramientas de reconocimiento automático de voz (ASR) el contenido de estas clases. 

A su vez, se presenta la arquitectura final del sistema junto con las alternativas consideradas y los motivos por los que se dejaron de lado. También se expone la experimentación asociada a cada etapa del diseño del asistente. 


\hfill \break
\keywords{RAG, LLM, OpenFing, ASR, Educación}

%%%%%%%%%%%%%%%%%%%%%%%%%%%%%%%%%%%%%%%%%%%%%%%%%%%%%%%%%%%%%%%%%%%%%%%%%%%%%%%%%%%%%%%%%%%%%%%%

